\section{RESEARCH DESIGN AND METHODOLOGY}\label{sec3}

\subsection{The Natural Experiment: US--China Tariff Escalation (2018--2019)}

Our identification strategy exploits the United States' staggered tariff escalation on Chinese imports between July 2018 and September 2019 as an external trade shock. The USTR implemented four major rounds: List 1 (\$34 billion; effective July 6, 2018), List 2 (\$16 billion; August 23, 2018), List 3 (\$200 billion; 10\% from September 24, 2018 and 25\% from May 2019), and List 4A (\$112 billion; 7.5\% from September 1, 2019). These policy actions created abrupt changes in trade costs and sourcing incentives across product categories.

The core empirical idea is that countries with stronger pre-shock exposure to the US--China trade nexus faced a larger indirect shock when bilateral tariffs escalated. We therefore use a difference-in-differences design with heterogeneous treatment intensity based on pre-shock exposure, and we test whether national productive complexity conditions countries' resilience responses.

Our main paper-period treatment timing follows the conceptual onset window used in the manuscript design: $\textit{Post}_{t}=1$ for 2019--2022, and 0 otherwise. We additionally run identification diagnostics and alternative timing checks (including placebo pseudo-shocks and event-time pre-trend tests) as robustness layers.

\subsection{Sample Construction and Data Sources}

The project uses a two-layer data structure.

First, we construct a broad country-year databank used for variable engineering and coverage checks. This databank contains 231 countries and 2,995 country-year observations from 2012 to 2024. Atlas-based trade and complexity variables are available beyond 2022, while TiVA-based GVC variables end in 2022.

Second, confirmatory estimation uses the 2012--2022 panel and model-specific complete-case samples after merging all controls and fixed effects covariates. As a result, estimation samples are unbalanced and differ by dependent variable:
\begin{itemize}
\item H1 (GVC linkage stability): 679 observations, 76 countries, years 2013--2021.
\item H2 (export recovery): 1,729 observations, 180 countries, years 2012--2021.
\item H3 (partner diversification): 1,729 observations, 180 countries, years 2012--2021.
\end{itemize}

Data sources are as follows:
\begin{itemize}
\item \textbf{Productive complexity and trade structure:} Harvard Growth Lab Atlas of Economic Complexity country-year and bilateral files (ECI, COI, export values, bilateral partner flows).
\item \textbf{GVC measures:} OECD Trade in Value Added (TiVA, 2023 release), used to construct forward-linkage metrics.
\item \textbf{Macroeconomic and policy controls:} World Development Indicators (WDI), including GDP per capita, trade openness, natural resource rents, and applied weighted-mean tariffs.
\item \textbf{Institutional controls:} Worldwide Governance Indicators (WGI), aggregated as a six-dimension governance composite.
\item \textbf{Geopolitical risk:} Caldara--Iacoviello monthly GPR series aggregated to annual values. We use country-specific annual GPR where available and otherwise use the global annual index as fallback for model completeness.
\end{itemize}

\subsection{Variable Operationalization}

\subsubsection{Dependent Variables}

\textbf{DV1: GVC linkage stability.} We begin with annual change in forward linkage, $\Delta \text{FEXGR\_DVA}_{it}$, from TiVA. For confirmatory estimation, we use a stability transformation:
\[
\text{GVC\_Stability}_{it} = -\left| \Delta \text{FEXGR\_DVA}_{it} \right|,
\]
so values closer to zero represent smaller disruption (higher stability).

\textbf{DV2: Export Recovery Index (ERI).} For each country-year, ERI is total exports divided by the country-specific pre-shock baseline average over 2015--2017:
\[
\text{ERI}_{it} = \frac{\text{Exports}_{it}}{\overline{\text{Exports}}_{i,2015\text{--}2017}}.
\]

\textbf{DV3: Trade partner diversification.} We compute destination concentration using the export-share HHI and define diversification as:
\[
\text{Diversification}_{it} = 1-\text{HHI}_{it}.
\]
Higher values indicate broader partner dispersion.

\subsubsection{Independent and Moderating Variables}

\textbf{Economic Complexity Index (ECI).} The focal explanatory variable is country-year ECI from Atlas. In regressions, ECI is standardized ($z$-score).

\textbf{Tariff shock exposure design.} We compute each country's pre-shock US--China trade intensity as the mean (2015--2017) of:
\[
\frac{\text{Trade with US and China}}{\text{Total trade}}.
\]
Countries in the top tercile are coded $\textit{Exposed}_i=1$ (otherwise 0). The treatment interaction is:
\[
\textit{PE}_{it}=\textit{Post}_{t}\times \textit{Exposed}_{i}.
\]
The key slope term for H1--H3 is $z\_ECI_{it}\times \textit{PE}_{it}$.

\textbf{H4 moderator (COI).} Product-space density is measured with standardized COI; the test term is $z\_ECI_{it}\times \textit{PE}_{it}\times z\_COI_{it}$.

\textbf{H5 moderator (GPR).} Geopolitical risk is measured with standardized annual GPR; the test term is $z\_ECI_{it}\times \textit{PE}_{it}\times z\_GPR_{it}$.

\subsubsection{Control Variables}

Baseline confirmatory specifications include standardized controls for:
\begin{itemize}
\item log GDP per capita,
\item trade openness (\% GDP),
\item WGI governance composite,
\item natural resource rents (\% GDP),
\item pre-shock US--China trade intensity.
\end{itemize}

Applied weighted-mean tariffs are available in the databank and in dedicated stricter-control builds, but are excluded from the main confirmatory specification to preserve cross-country coverage.

\subsection{Empirical Specification and Inference}

For each resilience outcome $Y_{it}$, we estimate:
\[
Y_{it}=\beta_1 zECI_{it}+\beta_2 PE_{it}+\beta_3 (zECI_{it}\times PE_{it})+\Gamma'X_{it}+\alpha_i+\lambda_t+\varepsilon_{it},
\]
where $\alpha_i$ are country fixed effects, $\lambda_t$ are year fixed effects, and standard errors are clustered at the country level. The parameter of interest for H1--H3 is $\beta_3$.

Moderation tests for H4 and H5 extend the equation with triple interactions:
\[
zECI_{it}\times PE_{it}\times zCOI_{it}
\quad\text{and}\quad
zECI_{it}\times PE_{it}\times zGPR_{it}.
\]

Inference is based on:
\begin{itemize}
\item cluster-robust country-level standard errors,
\item wild-cluster bootstrap p-values for primary H1--H3 terms,
\item Benjamini--Hochberg FDR adjustment across the confirmatory family.
\end{itemize}

\subsection{Identification Diagnostics and Transparency Protocol}

We report event-time pre-trend tests and placebo pseudo-shock tests as explicit identification diagnostics. In our confirmatory run, pre-trend tests do not reject parallel pre-trends for H1--H3, while the H1 placebo term is statistically significant; therefore, H1 is treated as not confirmed in strict confirmatory inference. H2 and H3 do not show significant focal effects under confirmatory criteria.

To avoid HARKing, global H1--H5 remain the primary confirmatory hypotheses. Any subgroup, regime, or income-class patterns are reported as exploratory heterogeneity analyses and interpreted as hypothesis-generating rather than confirmatory evidence.
